\chapter{Statistics and derived data}\label{ch:stats}
\begin{goals}
\textbf{In this chapter} you read the statistics card (the mean, its concentration and confidence, the
eigen analysis and the fabric shape), and keep results as datasets: the mean, the eigenvectors, a
$\beta$ diagram, and the line where two sets of planes meet.
\end{goals}

The statistics answer the first questions you ask of a set of orientations: where it points, how
tightly, and whether it is a cluster or a girdle. The \ui{statistics} tab\index{statistics tab}, under the
diagrams beside \ui{table} and \ui{console}, describes each visible layer of the active net: one card per
layer, in the layer's color. Show or hide a layer, edit its rows, or switch to another net, and the cards
follow.

\shot[0.8\linewidth]{stats-card}{The statistics card for the Tocher sample: an axial mean that means
little, a strong girdle (G~0.63, K~0.21), and its axis V3.}

\section{The card}\label{sec:stats-card}

The heading names the net and how many layers it shows (\ui{Net 1 · 1 layer}). Each card heads with the
layer's name and \ui{n}, the number of measurements wuffle could read. When some rows couldn't be read it
says so, as \ui{n 190 of 200}. Then, for the Tocher sample:

\begin{center}\small
\begin{tabularx}{\linewidth}{@{}>{\raggedright}p{0.17\linewidth}>{\raggedright}p{0.3\linewidth}L@{}}\toprule
\sffamily\bfseries line & \sffamily\bfseries Tocher & \sffamily\bfseries what it is \\\midrule
\ui{mean} & \code{027/89 · pole 207→01} & the axial Fisher mean; for planes, the mean plane and its pole \\
\ui{κ} & \code{3.1} & Fisher's concentration: larger is tighter about the mean \\
\ui{α95} & \code{6.8°} & the half-angle of the 95\,\% cone of confidence about the mean \\
\ui{S1 S2 S3} & \code{0.548 0.383 0.069} & the eigenvalues of the orientation tensor, scaled to sum to 1 \\
\ui{V1}, \ui{V2}, \ui{V3} & \code{205→01}, \code{295→42}, \code{113→48} & the eigenvectors, as lines \\
& \code{girdle 293/42} & for planes, the plane whose pole is V3: the best-fit great circle \\
\ui{K · C} & \code{0.21 · 2.07} & Woodcock's shape and strength \\
\ui{P G R} & \code{0.16 0.63 0.21} & Vollmer's point, girdle and random fractions \\
\bottomrule\end{tabularx}
\end{center}

Hover over a line's label for a reminder of what it is. A mean needs two measurements and the eigen
analysis three; until then the card says \ui{needs ≥ 2}. A fault layer adds the principal stresses,
\ui{σ1 σ2 σ3} and \ui{φ} (\cref{ch:faults}).

\section{The mean}\label{sec:stats-mean}

The mean is \term[Fisher mean]{Fisher's mean} \parencite{fisher1953} of the poles (or of the lines), taken \emph{axially}: a pole that
points down and the same axis pointing up count as one direction. wuffle reports it as an attitude, a plane
\emph{and} its pole for planes, a line for lines.

\begin{description}
\item[\ui{κ}] is Fisher's concentration parameter, estimated from how closely the measurements gather.
Values of a few, like Tocher's 3.1, mean a loose spread; a tight cluster of joints runs to tens or
hundreds.
\item[\ui{α95}] is the half-angle of the cone that holds the true mean with 95\,\% confidence. It shrinks
as you add measurements, even when they are scattered: Tocher's 200 loose poles give 6.8°.
\end{description}

\begin{caution}
A mean is only a summary of a cluster. Tocher's poles lie on a girdle, not about a point, so the mean
plane \code{027/89} is an average of very different planes and describes none of them. Read the eigen
analysis first; the fabric plot (\cref{ch:roses}) shows the same thing at a glance.
\end{caution}

\section{The eigen analysis}\label{sec:stats-eigen}

wuffle builds the \term{orientation tensor} of the poles (or lines) and finds its eigenvalues and eigenvectors:

\begin{description}
\item[\ui{S1 S2 S3}] the eigenvalues, largest first, scaled to sum to 1. A tight cluster has S1 near 1; a
perfect girdle has S1 and S2 near ½ and S3 near 0; a random spread has all three near ⅓.
\item[\ui{V1 V2 V3}] the eigenvectors, as lines. V1 is the direction of greatest concentration. For a
girdle of poles, V3 is the girdle's pole: the axis about which the planes turn, a cylindrical fold's
axis. The card gives the girdle itself too, the plane whose pole is V3 (Tocher's \code{293/42}).
\item[\ui{K · C}] shape and strength after \textcite{woodcock1977}\index{Woodcock's K and C}: $K = \ln(S_1/S_2) / \ln(S_2/S_3)$ and
$C = \ln(S_1/S_3)$. $K > 1$ is a cluster and $K < 1$ a girdle; a larger $C$ is a stronger fabric. Tocher's
0.21 is a girdle, and its 2.07 a well-defined one.
\item[\ui{P G R}] the fractions of \textcite{vollmer1990}\index{Vollmer's P, G and R}: $P = S_1 - S_2$, $G = 2(S_2 - S_3)$ and $R = 3S_3$.
They sum to 1 and place the fabric in a triangle of point, girdle and random. Tocher's G of 0.63 says the
same as its K: mostly girdle.
\end{description}

\section{Results as datasets}\label{sec:derived}

\shot{derived-net}{The mean (a ring with its cone of confidence, near \code{207→01}) and the
eigenvectors kept as a dataset, drawn in blue on the contoured Tocher poles.}

Each card has buttons that keep a result as a \term{derived dataset}:

\begin{description}
\item[\ui{mean → dataset}] one row: the mean's attitude, with \col{kappa}, \col{alpha95} and \col{n}
beside it. A plane for planes, a line for lines.
\item[\ui{eigenvectors → dataset}] three lines, V1 to V3, each with its \col{axis} name and
\col{eigenvalue}.
\item[\ui{β → dataset}] for planes: the line of intersection of every pair of planes, the classic
$\beta$ diagram\index{beta diagram}. Its \col{a} and \col{b} columns give the row numbers of the two planes each line came
from. 200 planes make 19\,900 lines; wuffle refuses past 600 planes (179\,700 lines) and asks you to select
fewer first.
\item[\ui{∩ intersect with…}] for planes, when another layer of planes is drawn on the same net: choose
it, and wuffle keeps the line where the two layers' \emph{mean} planes meet, with the \col{angle} between
them. Two limbs of a fold give its axis this way.
\end{description}

A derived dataset is drawn on the active net at once and opens in the table. Making it is one undo step,
\ui{Undo derive}. It is an ordinary dataset from then on: style it, draw it on other nets, select it,
save it. What makes it derived is its lineage (\cref{sec:lineage}). The tree marks it with an arrow, and hovering it
says how it was made (\ui{made by mean from tocher}). The lineage is kept in the project:

\begin{codeblock}
$ cat "data/tocher · mean.json"
  ...
  "lineage": {
    "op": "mean",
    "params": { "axial": true, "confidence": 0.95 },
    "sources": [ { "id": "d1", "name": "tocher" } ],
    "at": "2026-11-26T14:33:05.163Z"
  },
\end{codeblock}

A derived dataset is a snapshot, not a live formula. Edit \code{tocher} afterwards and \code{tocher · mean}
keeps the mean it was made with; make it again for the new one.

\begin{note}
A derived dataset that is drawn gets a statistics card of its own, like any layer. The card of three
eigenvectors, or of a single mean, describes those few lines and means little (three perpendicular axes
read as perfectly random). Hide the layer, or ignore its card.
\end{note}

\section{An intersection}\label{sec:stats-intersect}

\shot{intersect}{Tocher split into two sets by dip direction, \code{east} and \code{west}, with their
means and cones, and the line where the mean planes meet (violet, near \code{123→33}).}

To intersect two sets, they have to be two datasets drawn on one net. Split one with a selection
(\cref{sec:select-dataset}), or open two files. The Tocher sample split at a dip direction of 180°
gives:

\begin{center}\small
\begin{tabularx}{\linewidth}{@{}lrlrrL@{}}\toprule
\sffamily\bfseries set & \sffamily\bfseries n & \sffamily\bfseries mean & \sffamily\bfseries κ & \sffamily\bfseries α95 & \sffamily\bfseries K \\\midrule
\code{east} & 151 & \code{051/63} & 3.1 & 7.9° & 0.12 \\
\code{west} & 49 & \code{199/70} & 7.8 & 7.7° & 1.46 \\
\bottomrule\end{tabularx}
\end{center}

and \ui{∩ intersect with… › west} on the \code{east} card keeps \code{east ∩ west}: the line
\code{123→33}, the mean planes 56° apart.

\section{From the console}\label{sec:stats-console}
\index{stats (command)}\index{mean (command)}\index{eigen (command)}\index{beta (command)}%
\index{intersect (command)}%

\begin{codeblock}
stats                   # the first visible layer's card, as text
stats tocher | grep mean
mean                    # mean → dataset, for the first visible layer
eigen tocher            # eigenvectors → dataset
beta                    # the beta diagram
intersect east west     # where two layers' mean planes meet
\end{codeblock}

\code{stats} prints three lines: the count, the mean with $\kappa$ and $\alpha_{95}$, and the eigenvalues
with K and C:

\begin{codeblock}
tocher · n 200
mean 027/89 · pole 207→01 · κ 3.1 · α95 6.8°
S 0.548 0.383 0.069 · K 0.21 · C 2.07
\end{codeblock}

\code{mean}, \code{eigen}, \code{beta} and \code{intersect} answer with the new dataset's name and size,
\ui{→ tocher · β (19900 rows), drawn on Net 1}, so a routine can make all of them in a few lines
(\cref{sec:routines}).

\begin{tryit}
\begin{enumerate}
\item Open the Tocher sample and click the \ui{statistics} tab. Read K and G: is this a cluster or a
girdle?
\item Press \ui{eigenvectors → dataset}. Three blue points appear on the net; V3 sits in the gap the girdle
of poles wraps around.
\item In the console, type \code{beta}. 19\,900 intersection lines are drawn. Type
\code{stats "tocher · β"}: their mean is \code{114→48}, a degree from V3. Two ways to a fold axis, one
answer.
\item Type \code{cat "data/tocher · β.json"} and find the lineage that says where the lines came from.
\end{enumerate}
\end{tryit}
